% Vietnamese translation for OI Public Library
% Source-File: IOI中国国家候选队论文/2008/Day1/7.方戈《浅析信息学竞赛中一类与物理有关的问题》/3D_Music.doc
\documentclass[11pt]{article}
\usepackage{oipl-vn}

\title{Hàng rào nước âm nhạc ba chiều}
\author{POJ Monthly 2007.07.08}
\date{}

\begin{document}
\maketitle

\origtitle{3-dimensional Musical Water-fence}
\sourcepath{IOI 2008 / Fang Ge / 3D\_Music.doc}

\section{Mô tả}

Một lần Farmer John đang giặt quần áo bẩn. Nhìn rất nhiều xô nước, ông bỗng
nảy ra ý tưởng: ``Ta có thể làm một hàng rào nước âm nhạc ba chiều như thế
này.''

Sau một thời gian rất lâu, Farmer John đẽo một khối gỗ thành một hàng rào
nước âm nhạc ba chiều hoàn chỉnh, viết tắt là \term{3dmwf}. Như hình trong
nguồn gốc minh họa, 3dmwf là một mảng gồm \(n\) hàng và \(m\) cột các ô vuông.
Ô ở hàng \(i\), cột \(j\) được chiếm trọn bởi một khối \(b_{ij}\) có chiều
cao \(h_{ij}\). Phía trên mỗi khối có một vòi nước có thể đổ nước xuống.

Hình nguồn cho ví dụ một mảng \(4\times 4\) với chiều cao các khối:
\[
\begin{array}{rrrr}
5 & 12 & 11 & 8\\
15 & 2 & 4 & 14\\
16 & 1 & 3 & 7\\
13 & 10 & 9 & 6
\end{array}
\]

Nếu Farmer John mở một vòi, nước sẽ rơi xuống khối gỗ bên dưới do trọng lực.
Khi nước chạm vào khối, một âm thanh du dương vang lên. Cao độ của âm thanh
được quyết định bởi chiều cao của khối tạo ra âm thanh đó. Farmer John không
thích hai khối khác nhau tạo cùng một cao độ, nên không có hai khối nào có
cùng chiều cao. Để đơn giản, giả sử chiều cao của mỗi khối nằm trong khoảng
từ \(1\) đến \(n\times m\).

Ký hiệu \(\operatorname{pour}(i,j,k)\) là thao tác đổ \(k\) đơn vị nước lên
khối \(b_{ij}\). Sau khi đổ nước lên 3dmwf, có thể một phần nước sẽ bị giữ
lại trong các hố lõm.

\section{Quy tắc dòng chảy}

Khi thực hiện \(\operatorname{pour}(i,j,k)\), đúng một trong hai trường hợp
sau xảy ra.

\subsection{Chảy xuống các khối thấp hơn}

Nếu trên \(b_{ij}\) chưa có nước và tồn tại một số khối kề cạnh có mực ngang
thấp hơn \(b_{ij}\), đặt \(\mathit{Sum}\) là số khối kề như vậy. Khi đó
\(\operatorname{pour}(i,j,k)\) tương đương với các thao tác
\[
  \operatorname{pour}(i',j',k/\mathit{Sum})
\]
cho từng khối kề thấp hơn \(b_{i'j'}\). Nghĩa là nước chảy đều sang tất cả
các khối kề thấp hơn.

Với mảng ví dụ ở trên, nguồn gốc đưa ra chuỗi biến đổi:
\[
\begin{aligned}
\operatorname{pour}(4,1,27)
&\Longleftrightarrow
  \operatorname{pour}(4,0,9),\
  \operatorname{pour}(5,1,9),\
  \operatorname{pour}(4,2,9)\\
&\Longleftrightarrow
  \text{lãng phí }18\text{ đơn vị nước},\
  \operatorname{pour}(5,2,3),\
  \operatorname{pour}(4,3,3),\
  \operatorname{pour}(3,2,3)\\
&\Longleftrightarrow
  \text{lãng phí }21\text{ đơn vị nước},\
  \operatorname{pour}(5,3,1),\
  \operatorname{pour}(4,4,1),\
  \operatorname{pour}(3,3,1),\
  \operatorname{pour}(3,2,3)\\
&\Longleftrightarrow
  \text{lãng phí }22\text{ đơn vị nước},\
  \operatorname{pour}(5,4,0.5),\
  \operatorname{pour}(4,5,0.5),\
  \operatorname{pour}(3,3,1),\
  \operatorname{pour}(3,2,3)\\
&\Longleftrightarrow
  \text{lãng phí }23\text{ đơn vị nước},\
  \operatorname{pour}(3,3,1),\
  \operatorname{pour}(3,2,3).
\end{aligned}
\]
Nước của hai thao tác cuối bị giữ lại trong hố lõm tạo bởi
\(b_{22},b_{23},b_{32},b_{33}\). Vì vậy tổng lượng nước lãng phí cuối cùng là
\(23\) đơn vị.

\subsection{Nâng mực nước trong hố lõm}

Nếu trên \(b_{ij}\) đã có nước, hoặc không có khối kề nào có mực ngang thấp
hơn \(b_{ij}\), thì \(k\) đơn vị nước mới sẽ hòa với nước đang có trên
\(b_{ij}\). Mực nước ngang của khối chứa nước tăng lên cho đến khi không thể
tăng nữa vì có một lối thoát khiến nước chảy ra khỏi hố lõm.

Tiếp tục ví dụ trên:
\begin{itemize}
  \item Với \(\operatorname{pour}(3,2,3)\), sau khi đổ \(1\) đơn vị nước, các
  mực ngang liên quan trở thành
  \[
  \begin{array}{rrrr}
  5 & 12 & 11 & 8\\
  15 & 2 & 4 & 14\\
  16 & 2 & 3 & 7\\
  13 & 10 & 9 & 6
  \end{array}
  \]
  và sau khi đổ tiếp \(2\) đơn vị nước:
  \[
  \begin{array}{rrrr}
  5 & 12 & 11 & 8\\
  15 & 3 & 4 & 14\\
  16 & 3 & 3 & 7\\
  13 & 10 & 9 & 6
  \end{array}
  \]
  \item Với \(\operatorname{pour}(3,3,1)\), sau khi đổ \(1\) đơn vị nước,
  mực nước trong hố trở thành \(3+\frac13\):
  \[
  \begin{array}{rrrr}
  5 & 12 & 11 & 8\\
  15 & 3+\frac13 & 4 & 14\\
  16 & 3+\frac13 & 3+\frac13 & 7\\
  13 & 10 & 9 & 6
  \end{array}
  \]
  \item Nếu tiếp tục đổ rất nhiều nước vào hố lõm, mực nước sẽ đạt
  \[
  \begin{array}{rrrr}
  5 & 12 & 11 & 8\\
  15 & 7 & 7 & 14\\
  16 & 7 & 7 & 7\\
  13 & 10 & 9 & 6
  \end{array}
  \]
  và nước thừa sẽ chảy ra qua \(b_{34}\).
\end{itemize}

\section{Yêu cầu}

Vào một ngày nắng, sau khi làm việc trên nông trại rộng lớn, Farmer John rất
mệt và muốn nghe một chút âm nhạc du dương. Ông soạn một bản nhạc dài cho
3dmwf. Bản nhạc là một danh sách các thao tác \(\operatorname{pour}(i,j,k)\)
được sắp theo một thứ tự cụ thể. Ông mời bạn bè Wang Dong, Guo Huayang, Chen
Danqi và He Yijun giúp vận hành 3dmwf.

Farmer John là một nông dân tiết kiệm. Ông không muốn lãng phí quá nhiều nước
khi chơi nhạc bằng 3dmwf. Hãy cho ông biết tổng thể tích nước bị lãng phí.

\section{Dữ liệu vào}

Dòng đầu tiên chứa hai số \(n\) và \(m\), lần lượt là số hàng và số cột.

\(n\) dòng tiếp theo, mỗi dòng chứa \(m\) số nguyên, là chiều cao của các
khối gỗ.

Dòng tiếp theo chứa một số nguyên \(t\), là độ dài của bản nhạc.

\(t\) dòng tiếp theo, mỗi dòng chứa ba số nguyên \(i,j,k\), biểu diễn thao
tác \(\operatorname{pour}(i,j,k)\). Thứ tự các thao tác trong dữ liệu vào
chính là thứ tự trong bản nhạc.

\section{Dữ liệu ra}

In thể tích nước lãng phí tổng cộng, với đúng \(2\) chữ số sau dấu thập phân.

\section{Ví dụ}

\paragraph{Ví dụ vào 1.}
\begin{verbatim}
4 4
5 12 11 8
15 2 4 14
16 1 3 7
13 10 9 6
1
4 1 27
\end{verbatim}

\paragraph{Ví dụ ra 1.}
\begin{verbatim}
23.00
\end{verbatim}

\paragraph{Ví dụ vào 2.}
\begin{verbatim}
4 4
5 12 11 8
15 2 4 14
16 1 3 7
13 10 9 6
2
4 1 27
2 2 100
\end{verbatim}

\paragraph{Ví dụ ra 2.}
\begin{verbatim}
109.00
\end{verbatim}

\section{Giới hạn}

\[
\begin{gathered}
1\le n\le 1000,\qquad 1\le m\le 1000,\\
0\le t\le 1\,000\,000,\qquad 0\le k\le 1\,000\,000.
\end{gathered}
\]

\section{Nguồn}

POJ Monthly 2007.07.08, Chen Qifeng.

\end{document}
